\section{Discussion and Threats to Validity}
\label{sec:discussion}

\subsection{Diagnostic Explainability vs. Threshold Calibration}
\label{sec:disc:explainability}

The Diagnostic Pathway successfully provides deterministic, role-specific quality attribution (Reliability Engineer vs. SRE vs. Architect), translating abstract topological risk into actionable review decisions---and Section~\ref{sec:results:detection} shows this attribution, not superior ranking, is where its value over a scalar centrality score lies. That value is not free, however: Section~\ref{sec:results:calibration} shows the catalog's fixed $k=1.5$ IQR fence produces near-chance component-level agreement ($\kappa \approx 0$) despite high recall, because $98.4\%$ of active findings are edge-scoped and implicate components only as endpoints. Tightening $k$ would trade recall for precision on the same finding mix, but would not by itself change the underlying scope imbalance---a catalog dominated by two edge-level detectors (\texttt{BOTTLENECK\_EDGE}, \texttt{QOS\_MISMATCH}) will always implicate components indirectly at scale. The mitigation this paper actually validates is architectural rather than purely statistical: absolute gating blocks on every one of the eight benchmark scenarios (Section~\ref{sec:cicd:semantics}), which is unusable in practice, while delta-aware gating---blocking only newly introduced findings relative to a merge-base---bounds the review surface to what a given change actually introduces, regardless of the base rate. Threshold recalibration (Section~\ref{sec:conclusion:future_work}) and delta-aware gating are therefore complementary, not substitutes: one narrows what fires, the other narrows what is enforced.

\subsection{The Verify-Before-Recommend Discipline}
\label{sec:disc:verify}

Counterfactual verification proves essential: an unverified engine would have applied 461 harmful or ineffective mutations ($29.0\%$). By filtering mutations against cascade noise, SaG-Prescribe ensures that recommended refactorings reliably harden system architecture. Notably, this discipline is what makes Section~\ref{sec:results:prescriptive}'s reallocation result trustworthy rather than merely plausible: the mechanistic argument for why anti-affinity reallocation should work (removing SPOF co-location) is exactly the kind of confident-sounding reasoning an open-loop recommender would have shipped regardless of whether it was true, and it is only because every candidate is independently simulated that the $77.7\%$ survival rate is a measurement rather than an assumption.

\subsection{Construct Validity: Substrate Adequacy}
\label{sec:disc:substrate}

Beyond the general dependence on a discrete-event cascade simulator, this paper surfaces a more specific construct-validity threat, promoted here to first-class status rather than left implicit: a composite score's sub-dimensions can be individually well- or ill-posed on a given graph projection, and scoring on an ill-posed projection silently degrades the composite without changing its formula. Section~\ref{sec:model:substrate} states this as a design-time property; Section~\ref{sec:results:attribution} measures its concrete instance---the Availability dimension carrying almost no variance on the application-layer projection used throughout this paper's headline detection evaluation, because that projection has essentially no articulation points to find. The mitigation is not a different formula for $A(v)$, but scoring Availability-sensitive review on a projection where cut structures can occur (the middleware or system layer), and reporting per-layer results rather than treating the application layer as universally representative. We flag this as a threat other RM-style composite scores on typed dependency graphs should check for explicitly, not only fix locally: any weighted sum over structurally heterogeneous sub-scores inherits this risk whenever one term is a discrete graph predicate.

\subsection{Comparison with Prior Self-Reported Results}
\label{sec:disc:prior_results}

Section~\ref{sec:results:attribution} reconciles this paper's $\rho = 0.485$ against our own prior published $\rho = 0.94$ \cite{yigit2025graph}. In discussion terms, the general lesson is broader than this one paper pair: a criticality score whose ground-truth oracle is itself a connectivity statistic on the same graph the score is computed from will tend to correlate well with that oracle almost by construction, and that correlation does not necessarily transfer to an oracle with independent causal mechanics (here, a discrete-event message-dropout and queue-exhaustion simulator). We consider the re-measurement a strength of this submission's methodology rather than a weakness to minimize: reporting a negative result against one's own prior work, with the mechanism identified rather than merely disclosed, is what makes the finding actionable for other researchers evaluating structural criticality scores against connectivity-derived oracles.

\subsection{Further Threats to Validity}
\label{sec:disc:threats}

\begin{itemize}
    \item \textbf{Construct Validity:} Both detection and verification utilize a discrete-event cascade simulator. We mitigate simulator dependency by grounding operators in established distributed systems dependability patterns (anti-affinity scheduling, topic isolation, QoS hardening), and by validating the composite against SPOF ground truth directly (Section~\ref{sec:results:attribution}) rather than relying on ranking correlation alone.
    \item \textbf{Internal Validity:} Decoupled in-memory repositories ensure no circular data leakage between diagnostic candidate generation and counterfactual simulation.
    \item \textbf{External Validity:} Scenarios are synthetic domain models; validation on harvested industrial cyber-physical codebases remains high-priority future work.
    \item \textbf{Conclusion Validity:} The Wilcoxon test underlying Section~\ref{sec:results:prescriptive}'s significance claim has $n=7$, the minimum for which a two-sided signed-rank test can reach $p<0.05$; the reported $p=0.0156$ is the best attainable significance at this sample size, not evidence of an especially large effect.
\end{itemize}
